\renewcommand{\Anno}{1997}

\renewcommand{\NumeroProblema}{2}

\renewcommand{\Testo}{ 
La Scuola Superiore S. Anna possiede un auditorium a forma di triangolo equilatero in cui pareti, pavimento e soffitto assorbono perfettamente i suoni.  Ai vertici del triangolo sono disposte tre casse acustiche le cui potenze (misurate in centinaia di Watt) sono $1$, $1$, $k$, rispettivamente.

Come è ben noto, l'intensità del suono percepito in un punto è direttamente proporzionale alla potenza della sorgente ed inversamente proporzionale al quadrato della distanza dalla sorgente.
	
Determinare i valori di $k$ per cui esiste un punto della stanza in cui il suono emesso da ogni cassa è sentito alla stessa intensità e, per tali valori di $k$, determinare il punto.
}

\renewcommand{\Risposta}{
Esiste il punto se e solo se $k\leq 3$.  In tal caso il punto cercato appartiene all'asse del segmento che ha per estremi le 2 casse da 1, e la sua distanza dal segmento è data da
$$\frac{-\sqrt{3}+\sqrt{k(4-k)}}{2(k-1)}.$$
}

\renewcommand{\Aiutino}{
Consideriamo solo le 2 casse da 100 watt.  In quali punti il suono da loro emesso viene percepito con la stessa intensità?  Come varia in tale luogo l'intensità percepita del suono delle 3 casse? Questo basta per dire quando esiste il punto richiesto.  Per individuarlo esattamente, occorre poi impostare un'equazione opportuna.
}

\renewcommand{\SoluzioneUfficiale}{
Indichiamo con $A$, $B$, $C$ i vertici del triangolo, che possiamo supporre di lato unitario in qualche opportuna unità di misura, e supponiamo che le 2 casse da 1 stiano nei vertici $B$ e $C$. Avendo la stessa potenza, il suono emesso da queste 2 casse si percepirà con uguale intensità nei punti equidistanti da $B$ e $C$, cioè nei punti dell'asse del segmento $BC$, la cui intersezione con il triangolo è la mediana $AM$.

Consideriamo ora un punto $P$ variabile su tale mediana, ed indichiamo con $x$ la sua distanza da $M$, come nella figura qui sotto.
\begin{center}
	\psset{unit=10mm}
	\pspicture(0,-0.4)(4,3.8)
	
	\pstGeonode[PosAngle={210,-30}](0,0){B}(4,0){C}
	\pstRotation[PosAngle=90]{B}{C}{A}
	\pstLineAB{A}{B}
	\pstLineAB{C}{B}
	\pstLineAB{A}{C}
	\pstMiddleAB[PosAngle=-90]{B}{C}{M}
	\pstLineAB[linecolor=blue]{A}{M}
	\pstGeonode(2,1){P}
	\pstLineAB[linecolor=blue,linestyle=dashed]{B}{P}
	\rput[r](1.9,0.5){$x$}
	\endpspicture
\end{center}

Indichiamo con $I_{A}(x)$, $I_{B}(x)$, $I_{C}(x)$, rispettivamente, le intensità con cui viene percepito in $P$ il suono proveniente dalle casse poste in $A$, $B$, $C$.  Tenendo conto che $BM=CM=1/2$, $AM=\sqrt{3}/2$, quindi $AP=\sqrt{3}/2-x$, e per il teorema di Pitagora
$$BP^{2}=CP^{2}=x^{2}+\frac{1}{4},$$
avremo che
$$I_{B}(x)=I_{C}(x)=\frac{\alpha}{\displaystyle{x^{2}+\frac{1}{4}}},
\hspace{4em}
I_{A}(x)=\frac{\alpha k}{\displaystyle{\left(
\frac{\sqrt{3}}{2}-x\right)^{2}}},$$
dove $\alpha$ è un'opportuna costante di proporzionalità.

Da queste formule segue immediatamente che l'intensità del suono proveniente da $B$ e $C$ diminuisce con continuità all'aumentare di $x$, mentre l'intensità del suono proveniente da $A$ cresce con continuità all'aumentare di $x$, e tende a $+\infty$ quando $P$ si avvicina ad $A$ (qui stiamo supponendo, per ovvie ragioni fisiche, che sia $k>0$).

Pertanto un punto in cui le 3 intensità coincidono esisterà (e sarà unico) se e solo se in $M$ il suono proveniente da $A$ giunge con intensità minore od uguale del suono proveniente da $B$ e $C$. Poiché $M$ corrisponde ad $x=0$, questa condizione equivale a
$$\frac{\alpha k}{\left(\sqrt{3}/2\right)^{2}}\leq \frac{\alpha}{1/4},
$$
da cui con semplici calcoli si ottiene la condizione $k\leq 3$, da aggiungere alla condizione fisica $k>0$.

Il grafico seguente riporta l'andamento di $I_{A}(x)$ e $I_{B}(x)$ per alcuni valori del parametro $k$ (abbiamo supposto $\alpha=1$).
\begin{center}
	\psset{xunit=25ex,yunit=5ex} 
	\pspicture(0,-1)(2,9.5) 
	
	\psplot[linewidth=1.5\pslinewidth]{0}{3 sqrt 2 div}{1 x 2 exp 1 4 div add div}
	\psplot[linewidth=1.5\pslinewidth,linecolor=red]{0}{3 sqrt 2 div 1 3 div sub}{1 3 sqrt 2 div x sub 2 exp div}
	\psplot[linewidth=1.5\pslinewidth,linecolor=blue]{0}{3 sqrt 2 div 2 sqrt 3 div sub}{2 3 sqrt 2 div x sub 2 exp div}
	\psplot[linewidth=1.5\pslinewidth,linecolor=green]{0}{3 sqrt 2 div 3 sqrt 3 div sub}{3 3 sqrt 2 div x sub 2 exp div}
	\psplot[linewidth=1.5\pslinewidth,linecolor=cyan]{0}{3 sqrt 2 div 4 sqrt 3 div sub}{4 3 sqrt 2 div x sub 2 exp div}
	\psplot[linewidth=1.5\pslinewidth,linecolor=yellow]{0}{3 sqrt 2 div 1 2 div sqrt 3 div sub}{1 2 div 3 sqrt 2 div x sub 2 exp div}
	
	\psaxes[linewidth=.6\pslinewidth]{->}(1,9.5)
	\rput(!3 sqrt 2 div 0){\pspicture(0,0)(0,0)\psline[linewidth=.6\pslinewidth](0,-3pt)(0,3pt)\endpspicture}
	\rput[t](!3 sqrt 2 div 0.3 neg){$\sqrt{3}/2$}
	\psline[linewidth=1.5\pslinewidth](1.3,7)(1.5,7)
	\rput[l](1.6,7){$I_{B}(x)=I_{C}(x)$}
	\psline[linewidth=1.5\pslinewidth,linecolor=cyan](1.3,6)(1.5,6)
	\rput[l](1.6,6){$I_{A}(x)$ con $k=4$}
	\psline[linewidth=1.5\pslinewidth,linecolor=green](1.3,5)(1.5,5)
	\rput[l](1.6,5){$I_{A}(x)$ con $k=3$}
	\psline[linewidth=1.5\pslinewidth,linecolor=blue](1.3,4)(1.5,4)
	\rput[l](1.6,4){$I_{A}(x)$ con $k=2$}
	\psline[linewidth=1.5\pslinewidth,linecolor=red](1.3,3)(1.5,3)
	\rput[l](1.6,3){$I_{A}(x)$ con $k=1$}
	\psline[linewidth=1.5\pslinewidth,linecolor=yellow](1.3,2)(1.5,2)
	\rput[l](1.6,2){$I_{A}(x)$ con $k=1/2$}
	\endpspicture
\end{center}

Nel caso in cui $0<k\leq 3$, si può trovare il valore esatto di $x$ risolvendo l'equazione $I_{B}(x)=I_{C}(x)=I_{A}(x)$, e cioè
$$\frac{\alpha k}{\displaystyle{\left(
\frac{\sqrt{3}}{2}-x\right)^{2}}}=
\frac{\alpha}{\displaystyle{x^{2}+\frac{1}{4}}},$$
la quale con qualche calcolo si trasforma nell'equazione
$$(k-1)x^{2}+\sqrt{3}\,x+\frac{k-3}{4}=0.$$

Se $k=1$, tale equazione è di primo grado ed ha l'unica soluzione $x=\sqrt{3}/6$ che corrisponde, come potevamo aspettarci, al centro del triangolo.

Se $k\neq 1$, tale equazione è di secondo grado ed ha le soluzioni
$$x_{1}=\frac{-\sqrt{3}-\sqrt{k(4-k)}}{2(k-1)},
\hspace{6em}
x_{2}=\frac{-\sqrt{3}+\sqrt{k(4-k)}}{2(k-1)}.$$

Per quanto detto sopra, noi già sappiamo che per $0<k\leq 3$ una ed una sola di tali soluzioni appartiene all'intervallo $[0,\sqrt{3}/2)$, dunque rappresenta un punto $P$ accettabile. Dimostriamo che tale soluzione è $x_{2}$.

Infatti, moltiplicando numeratore e denominatore per $\sqrt{3}+\sqrt{k(4-k)}$, possiamo riscrivere $x_{2}$ nella forma
$$x_{2}=\frac{-k^{2}+4k-3}{2(k-1)(\sqrt{3}+\sqrt{k(4-k)})}=
\frac{3-k}{2(\sqrt{3}+\sqrt{k(4-k)})}.$$

Ora è evidente che $x_{2}\geq 0$ per $k\leq 3$. Per dimostrare che $x_{2}$ è anche minore di $\sqrt{3}/2$, dobbiamo risolvere la disequazione
$$\frac{3-k}{2(\sqrt{3}+\sqrt{k(4-k)})}<\frac{\sqrt{3}}{2}.$$

Moltiplicando per il denominatore (sicuramente positivo) e svolgendo i calcoli, si ritrova la disequazione
$$-k< \sqrt{3k(4-k)},$$
banalmente verificata per tutti i valori $0<k\leq 3$.
}

\renewcommand{\SoluzioniAlternative}{
}

\renewcommand{\Osservazioni}{\TitoloOsservazione
Nella soluzione presentata abbiamo ottenuto la condizione su $k$ ragionando sulla crescenza e decrescenza delle intensità percepite. Avremmo anche potuto limitarci ad impostare l'equazione $I_{A}(x)=I_{B}(x)=I_{C}(x)$ e, \emph{dopo} aver ottenuto le soluzioni $x_{1}$ e $x_{2}$, discutere al variare di $k$ la loro appartenenza all'intervallo di accettabilità $[0,\sqrt{3}/2)$. Questo però ci avrebbe portato a discutere ben 4 disequazioni irrazionali.

Per i lettori interessati a questo approccio, enunciamo le conclusioni di tale discussione senza riportarne i passaggi: ferme restando le condizioni $k\neq 1$ (altrimenti l'equazione è di primo grado e va studiata a parte) e $k\in[0,4]$ (altrimenti il discriminante è negativo), abbiamo che $x_{1}<0$ per ogni $k>1$ e $x_{1}>\sqrt{3}/2$ per ogni $k<1$ (quindi in particolare $x_{1}$ non sta mai nell'intervallo considerato), mentre $x_{2}\geq 0$ per $k\leq 3$ e
$x_{2}<\sqrt{3}/2$ per $k>0$.  }

\renewcommand{\Commenti}{
}